% mazzi: 20,21
\chapter{L'apparato respiratorio}

% ---------------------------------------------------------------------------
\section{Generalità}

\textbf{Sistema respiratorio}: assicura gli \textbf{scambi gassosi} con l'ambiente esterno
$\rightarrow$ l'organismo assume l'$O_2$ dell'aria e si libera della $CO_2$ prodotta dal metabolismo
cellulare.

\begin{table}[htbp]
  \centering \small \renewcommand{\arraystretch}{1.15}
  \begin{tabular}{|l|r|}
    \hline
    \textbf{Componente} & \textbf{Concentrazione} \\
    \hline
    Azoto ($N_2$)              & 78,084\% \\
    Ossigeno ($O_2$)          & 20,946\% \\
    Argon (Ar)                & 0,934\% \\
    Anidride carbonica ($CO_2$) & 360 ppmv \\
    Neon (Ne)                 & 18,18 ppmv \\
    Elio (He)                 & 5,24 ppmv \\
    Metano ($CH_4$)           & 1,6 ppmv \\
    Kripto (Kr)               & 1,14 ppmv \\
    Idrogeno ($H_2$)          & 0,5 ppmv \\
    Protossido di azoto ($N_2O$) & 0,3 ppmv \\
    Xenon (Xe)                & 0,087 ppmv \\
    \hline
  \end{tabular}
\end{table}

\textbf{Costituzione}: organi cavi (le \textbf{vie aeree}) + due organi parenchimatosi (i
\textbf{polmoni}).
\begin{itemize}
  \item \textbf{vie aeree}: canali che portano l'aria dall'ambiente esterno ai polmoni;
  \item \textbf{polmoni}: sede degli scambi gassosi tra aria e sangue ($O_2$ al sangue, $CO_2$
    eliminata);
  \item la \textbf{respirazione} consta di due fasi: \textbf{inspirazione} (l'aria con $O_2$ giunge
    ai polmoni) ed \textbf{espirazione} (l'aria con $CO_2$ esce all'esterno);
  \item funzioni accessorie: \textbf{olfattiva} (naso) e \textbf{fonazione} (organo preposto = la
    laringe).
\end{itemize}

\rubrica{Organi per regione}
\begin{itemize}
  \item \textbf{testa}: cavità nasale (passaggio e ripulitura dell'aria, olfatto, fonazione), seni
    paranasali (umidificazione/riscaldamento, alleggerimento del massiccio facciale);
  \item \textbf{collo}: rinofaringe e laringe (passaggio dell'aria, fonazione), trachea (passaggio,
    rimozione del pulviscolo, umidificazione/riscaldamento);
  \item \textbf{torace}: bronchi e polmoni (scambi gassosi tra aria e sangue; il polmone anche
    equilibrio acido-base), pleure (creano una cavità a pressione negativa per la dilatazione dei
    polmoni).
\end{itemize}
Le vie aeree superiori continuano, attraverso la faringe, nelle vie aeree inferiori (laringe,
trachea, bronchi). I bronchi principali dx e sx assicurano la ventilazione dei due polmoni.

\rubrica{Riassunto delle funzioni}
\begin{itemize}
  \item scambi gassosi tra aria e sangue circolante;
  \item protezione delle superfici respiratorie (disidratazione, sbalzi di temperatura, variazioni
    ambientali);
  \item difesa da microrganismi patogeni;
  \item produzione dei suoni;
  \item regolazione di volume e pressione del sangue, controllo del pH e dei fluidi corporei.
\end{itemize}

% ---------------------------------------------------------------------------
\section{Il naso e le cavità paranasali}

\textbf{Naso}: costituito in parte dalle ossa nasali e dalla mascella, in parte dalla cartilagine
del setto e dalle cartilagini alari (maggiore e minore). Funzioni:
\begin{enumerate}
  \item respirazione;
  \item climatizzazione dell'aria;
  \item percezione olfattiva;
  \item modulazione della voce.
\end{enumerate}
$\sim$6 l di aria al minuto a riposo. La \textbf{forma del naso} è un carattere etnico (in base al
rapporto altezza/larghezza dell'apertura piriforme: stretto e affilato, basso e schiacciato, o
intermedio). È la 1\textsuperscript{a} parte delle vie aeree superiori $\rightarrow$ responsabile di
riscaldamento, umidificazione e in parte filtrazione dell'aria inspirata.

\rubrica{Naso esterno}
Rilievo impari e mediano, a piramide triangolare con base aderente alla faccia.
\begin{itemize}
  \item 3 facce: 2 laterali e 1 inferiore; la faccia inferiore comunica con l'esterno tramite due
    orifizi, le \textbf{narici};
  \item rivestito da cute sottile e mobile; sistema muscolare = muscoli mimici (dilatatori/costrittori).
\end{itemize}

\rubrica{Seni paranasali}
\begin{itemize}
  \item cavità nasali delimitate da cartilagini fibro-elastiche e ossa;
  \item internamente cavità pneumatiche rivestite da mucosa respiratoria = \textbf{seni paranasali};
  \item seni paranasali e condotti nasolacrimali drenano nella cavità nasale (orifizi sulle pareti
    laterali);
  \item comprendono:
    \begin{enumerate}
      \item \textbf{seno frontale}: scavato nell'osso frontale;
      \item \textbf{labirinto etmoidale}: celle etmoidali (variabili per numero e ampiezza),
        nell'osso etmoide $\rightarrow$ le celle anteriori drenano presso gli osti di seno mascellare
        e frontale, le posteriori nel meato superiore (recesso sfeno-etmoidale);
      \item \textbf{seno sfenoidale}: nel corpo dell'osso sfenoide;
      \item \textbf{seno mascellare}: nel corpo della mascella.
    \end{enumerate}
\end{itemize}

% ---------------------------------------------------------------------------
\section{La faringe}

\textbf{Faringe}: organo cavo, impari e mediano, dalla base del cranio fino a C6 (dove continua con
l'esofago); anteriore alla colonna, posteriore a cavità nasali, cavità orale e laringe. È il
\textbf{crocevia} che mette in comunicazione con laringe ed esofago. Tre segmenti:
\begin{itemize}
  \item \textbf{rinofaringe} (superiore);
  \item \textbf{orofaringe} (medio);
  \item \textbf{laringofaringe} (inferiore).
\end{itemize}

\rubrica{Rinofaringe}
\begin{itemize}
  \item segmento connesso funzionalmente al solo sistema respiratorio; dalla base del cranio al
    palato molle (velo palatino);
  \item rivestita da mucosa respiratoria con accumuli di tessuto linfoide $\rightarrow$ tonsille
    faringea e tubariche.
\end{itemize}

% ---------------------------------------------------------------------------
\section{La laringe}

\textbf{Laringe}: organo cavo, impari e mediano, $\sim$4 cm; nel collo, dopo la faringe, dietro la
lingua, continua nella trachea. A piramide (base superiore, apice tronco in continuità con la
trachea); proiezione sulla colonna tra C4 e C6.
\begin{itemize}
  \item funzione: lasciarsi attraversare dall'aria in inspirazione ed espirazione; connessione
    anatomica e funzionale con i polmoni;
  \item \textbf{legamenti vestibolari e vocali} rivestiti da pieghe mucose proiettate verso la
    glottide; i legamenti vestibolari (nelle pieghe vestibolari) prevengono l'entrata di corpi
    estranei nella glottide e proteggono le pieghe vocali.
\end{itemize}

\rubrica{Fonazione}
L'aria che passa attraverso la glottide fa vibrare le pieghe vocali $\rightarrow$ onde sonore
(\textbf{fonazione}). L'intensità del suono dipende da diametro, lunghezza e tensione delle pieghe
vocali. Patologie: faringite, fratture da sport (es.\ hockey), incidenti stradali, tumori.

\rubrica{Laringe durante la deglutizione}
I muscoli intrinseci ed estrinseci associati alla laringe cooperano per evitare che cibo o liquidi
entrino nella glottide.

% ---------------------------------------------------------------------------
\section{La trachea}

\textbf{Trachea}: condotto aereo che fa seguito alla laringe; $\sim$11 cm, dalla cartilagine
cricoidea alla biforcazione. A \textbf{T4} si biforca in 2 \textbf{bronchi principali} (uno per
polmone).
\begin{itemize}
  \item 15--20 anelli cartilaginei incompleti, connessi da \textbf{legamenti anulari} connettivali;
  \item parete posteriore priva di scheletro cartilagineo, piatta, in rapporto con l'esofago
    $\rightarrow$ ne permette la distensione durante la deglutizione;
  \item il rilasciamento del muscolo tracheale è dovuto all'attivazione del \textbf{simpatico}
    $\rightarrow$ aumento del diametro tracheale.
\end{itemize}
Clinica correlata: riflesso della tosse, broncocostrizione, manovra di Heimlich, tracheotomia.

% ---------------------------------------------------------------------------
\section{I bronchi e l'albero bronchiale}

\textbf{Bronchi principali} (dx e sx): canali aeriferi che originano dalla biforcazione della
trachea, diretti in basso e lateralmente fino all'ilo del rispettivo polmone.
\begin{itemize}
  \item lume pervio grazie agli anelli cartilaginei; bronco \textbf{dx}: 4--6 anelli, bronco
    \textbf{sx}: 9--12;
  \item il bronco dx ha diametro maggiore del sx;
  \item mucosa ciliata e ricca di \textbf{cellule mucipare caliciformi} $\rightarrow$ catturano il
    pulviscolo e lo spingono verso l'esterno col movimento ciliare; l'occlusione porta a morte per
    asfissia;
  \item la componente cartilaginea si riduce progressivamente; radici dei polmoni anteriormente a
    T5 (dx) e T6 (sx).
\end{itemize}

\rubrica{Albero bronchiale}
\begin{enumerate}
  \item trachea;
  \item bronchi principali (dx e sx);
  \item bronchi lobari (I ordine: 3 nel polmone dx, 2 nel sx);
  \item bronchi segmentali o zonali (II ordine);
  \item bronchi interlobulari (III ordine);
  \item bronchi lobulari;
  \item bronchi intralobulari;
  \item bronchioli terminali $\rightarrow$ ventilano il \textbf{lobulo};
  \item bronchioli respiratori $\rightarrow$ organizzano l'\textbf{acino} (unità funzionale);
  \item dotti alveolari;
  \item sacco alveolare;
  \item alveoli.
\end{enumerate}
Differenza istologica: i \textbf{bronchi} presentano cartilagine nelle pareti (la ialina è
progressivamente sostituita da cartilagine elastica); i \textbf{bronchioli} non hanno cartilagine,
solo muscolatura liscia.

% ---------------------------------------------------------------------------
\section{I polmoni}

\textbf{Polmoni} (dx e sx): organi pari nella cavità toracica, a cono tronco con punta fino alla base
del collo (sopra la I costa).
\begin{itemize}
  \item base sulla faccia superiore del diaframma; apice arrotondato;
  \item facce: \textbf{diaframmatica} (concava, sulla cupola diaframmatica), \textbf{costo-vertebrale}
    (convessa, si adatta alla parete toracica), \textbf{mediastinica} (concava);
  \item al centro della faccia mediastinica: l'\textbf{ilo del polmone};
  \item \textbf{polmone dx}: 2 scissure $\rightarrow$ 3 lobi; \textbf{sx}: 1 scissura $\rightarrow$ 2
    lobi; il dx è più voluminoso del sx;
  \item consistenza spugnosa ed elastica; colore roseo nel bambino, scuro nell'adulto (pulviscolo
    inspirato e fagocitato dai macrofagi).
\end{itemize}

\rubrica{Zone e lobuli}
\begin{itemize}
  \item dopo il lobo, l'unità costitutiva indipendente è la \textbf{zona} (o segmento): territorio
    del lobo con indipendenza anatomica; ogni bronco lobare dà 2--5 bronchi zonali $\rightarrow$
    dx \textbf{10 zone}, sx \textbf{9 zone};
  \item ogni zona $\rightarrow$ centinaia di \textbf{lobuli}; ogni lobulo = un bronchiolo terminale +
    il tessuto polmonare associato;
  \item \textbf{acino}: insieme di ramificazioni con alveoli che originano da un bronchiolo
    respiratorio $\rightarrow$ \textbf{unità funzionale elementare} del parenchima.
\end{itemize}

% ---------------------------------------------------------------------------
\section{Le pleure e la cavità pleurica}

\textbf{Pleura}: membrana sierosa che avvolge ogni polmone, con due foglietti; bagnata da un liquido
giallastro filtrato dal plasma $\rightarrow$ facilita i movimenti dei polmoni.
\begin{itemize}
  \item \textbf{foglietto parietale}: aderisce alla fascia endo-toracica;
  \item \textbf{foglietto viscerale}: aderisce alla superficie del polmone, spingendosi nelle
    scissure;
  \item i due foglietti si riflettono l'uno nell'altro a livello dell'ilo;
  \item consentono lo scorrimento dei polmoni; la \textbf{pressione negativa} della cavità permette
    l'espansione in inspirazione e la compressione in espirazione.
\end{itemize}

\textbf{Cavità pleurica}: spazio toracico delimitato dal foglietto parietale. Insieme al
\textbf{mediastino} (cuore, timo, vasi, nervi, esofago) occupa l'intero torace.

% ---------------------------------------------------------------------------
\section{Gli alveoli}

\textbf{Alveoli}: parete di epitelio pavimentoso semplice, con due tipi cellulari + macrofagi.

\rubrica{Pneumociti di I tipo (cellule squamose alveolari)}
\begin{itemize}
  \item piatte, ampio citoplasma; formano la quasi totalità della parete ($\sim$97\% della superficie
    alveolare);
  \item organelli raggruppati attorno al nucleo; desmosomi e giunzioni occludenti;
  \item costituiscono la \textbf{barriera di diffusione dei gas}.
\end{itemize}

\rubrica{Pneumociti di II tipo}
\begin{itemize}
  \item corti microvilli; secernono il \textbf{surfactante} $\rightarrow$ abbassano la tensione
    superficiale degli alveoli.
\end{itemize}

\rubrica{Surfactante}
Emulsione lipidica secreta dai pneumociti di II tipo; riduce la tensione superficiale del fluido che
riveste l'alveolo $\rightarrow$ impedisce il collasso degli alveoli piccoli e l'espansione eccessiva
di quelli grandi (contrasta la chiusura degli alveoli, \textbf{legge di Laplace}). Composizione: 94\%
lipidi, 6\% proteine.
\begin{table}[htbp]
  \centering \small \renewcommand{\arraystretch}{1.15}
  \begin{tabular}{|l|r|}
    \hline
    \textbf{Componente} & \textbf{\%} \\
    \hline
    Fosfatidilcolina satura   & 50 \\
    Fosfatidilcolina insatura & 20 \\
    Fosfatidilglicerolo       & 8 \\
    Altri fosfolipidi         & 8 \\
    Lipidi neutri             & 8 \\
    Proteine                  & 6 \\
    \hline
  \end{tabular}
\end{table}

\rubrica{Macrofagi alveolari}
\begin{itemize}
  \item ``cellule della polvere'', derivano dai monociti;
  \item nella parete alveolare o negli alveoli stessi;
  \item fagocitano le piccole particelle sfuggite ai sistemi di protezione delle alte vie respiratorie.
\end{itemize}

\rubrica{Organizzazione alveolare}
Le fibre elastiche mantengono la posizione di alveoli e bronchioli respiratori e facilitano
l'espirazione promuovendo la riduzione delle dimensioni degli alveoli.

% ---------------------------------------------------------------------------
\section{La meccanica respiratoria}

\textbf{Respirazione}: azione coordinata di diaframma, muscoli addominali, toracici e cervicali.
\textbf{Obiettivo}: aumentare (inspirazione) o ridurre (espirazione) i diametri sagittale, verticale
e trasverso della cavità toracica.

\rubrica{Dinamica respiratoria}
\begin{itemize}
  \item a riposo, la funzione è legata prevalentemente alle escursioni del \textbf{diaframma};
  \item \textbf{inspirazione}: escursione verso il basso del centro tendineo del diaframma (favorita
    dal rilassamento addominale) $\rightarrow$ sollevamento delle coste e ampliamento della gabbia
    toracica;
  \item \textbf{espirazione}: intercostali interni + diaframma; se \textbf{forzata}, anche i muscoli
    dell'addome (retto, obliquo esterno/interno, trasverso).
\end{itemize}

\rubrica{Diaframma}
Muscolo impari, cupoliforme e laminare, che separa la cavità toracica da quella addominale.
\begin{itemize}
  \item spessore $\sim$3 mm; innervato dal \textbf{nervo frenico};
  \item a riposo rilassato; si contrae $\rightarrow$ il volume toracico aumenta; si rilassa
    $\rightarrow$ il volume diminuisce.
\end{itemize}

% ---------------------------------------------------------------------------
\section{Il controllo nervoso della respirazione}

\textbf{Centri respiratori}: il ritmo basale e la profondità del respiro sono controllati dal
\textbf{centro del ritmo del respiro} (gruppi respiratori \textbf{ventrale} e \textbf{dorsale}).
\begin{itemize}
  \item si trovano nella formazione reticolare, connessi con quasi tutti i nuclei sensitivi e motori
    $\rightarrow$ gli stati emotivi possono modificare ritmo e profondità del respiro;
  \item i \textbf{centri superiori} influenzano la respirazione inviando segnali al centro
    pneumotassico o agendo direttamente sui muscoli respiratori.
\end{itemize}

\rubrica{Bulbo (midollo allungato)}
Contiene i \textbf{centri del ritmo respiratorio}: controllano la frequenza di base dei movimenti
respiratori, a loro volta regolata dagli impulsi dei centri \textbf{apneustico} e
\textbf{pneumotassico} del ponte.

\rubrica{Ponte}
La formazione reticolare del ponte (su ogni lato dell'encefalo) contiene centri respiratori che
modificano l'attività dei centri del ritmo del respiro del midollo allungato.

% ---------------------------------------------------------------------------
\section{La circolazione polmonare}

\textbf{Circolazione polmonare} (piccola circolazione): circuito di vasi dal ventricolo dx
$\rightarrow$ si capillarizza a livello degli alveoli $\rightarrow$ torna all'atrio sx tramite le
vene polmonari (che trasportano sangue ossigenato).
\begin{itemize}
  \item sistema ad \textbf{elevato flusso} e \textbf{bassa pressione};
  \item la pressione idrostatica deve garantire che la pressione colloidosmotica renda il flusso in
    entrata superiore a quello in uscita;
  \item se ciò non accade $\rightarrow$ \textbf{edema polmonare} (accumulo di liquido negli alveoli:
    ipertensione, patologie renali/cardiache, infezioni).
\end{itemize}

% ---------------------------------------------------------------------------
\section{Note cliniche}

\rubrica{Effetti del fumo}
\begin{enumerate}
  \item danno e distruzione della parete degli alveoli;
  \item danno tossico alle cellule ciliate;
  \item riduzione e arresto del moto ciliare (morte delle cellule ciliate);
  \item maggiore suscettibilità alle infezioni;
  \item stato infiammatorio cronico.
\end{enumerate}

\rubrica{Altre patologie}
\begin{itemize}
  \item \textbf{mesotelioma} (tumore della pleura);
  \item \textbf{edema polmonare da alta quota}: accumulo di liquido negli alveoli per la rarefazione
    della pressione atmosferica in quota;
  \item \textbf{TBC} (tubercolosi);
  \item \textbf{COVID-19}: malattia virale (SARS-CoV-2) che può evolvere in \textbf{polmonite
    interstiziale} grave, con complicanze (trombosi polmonare, miocardite, nefropatia); nel modello
    proposto dalla docente, l'esercizio fisico intenso può favorire l'apertura degli alveoli e la
    penetrazione del virus, mentre una risposta immunitaria eccessiva provoca la ``tempesta
    citochinica''.
\end{itemize}
